% displayname : Ondes électromagnétiques dans le vide
\documentclass[12pt,fancy]{cours}
\Chapit{6}{3}
\begin{document}
\maketitre[***]{
Nous allons étudier la propagation de l'onde électromagnétique dans le vide illimité (absence de sources $\rho=0$ et $\Vec{j}=0$, et de conditions aux limites). Nous allons là encore retrouver un régime non dispersif, où la propagation de l'onde électromagnétique est décrite par une équation de d'Alembert. Nous étudierons surtout les OPPH: les équations de Maxwell étant linéaires, par transformée de Fourier l'étude de l'OPPH est la brique élémentaire pour comprendre la propagation d'un signal quelconque.
}{}

\ProgrammeColle

%%%%%%%%%%%%%%%%%%%%%%%%%%%%%%%%%%%%%%%%%%%%%%%%%%%%%%%%%%%%%%%%%%%%%%%%%%
\section{Équation d'onde de d'Alembert}
Les équations de couplage pour les grandeurs duales ($\vv{E},\vv{B}$) sont données par les équations de Maxwell. Dans ce chapitre nous étudions le champ électromagnétique dans la situation suivante.
\begin{enumerate}[label=\boxed{\rm H\arabic*}]
  \item Le milieu est le \imp{vide}, qui en particulier ne contient ni charge, ni courant, soit 
\eqbox[]{
  \rho=0 \text{ et } \j=0.
  }
  \begin{remarque}
    
Les équations de Maxwell se simplifient alors de la sorte.
  \begin{liste}
  \item Les équations aux divergences (M.G. et M.T.) informent sur la structure des champs $\E$ et $\B$, 
  %  \begin{equation}
  \eqbox[]{
    \begin{cases}
      \dive{\vv{E}}=0 & \text{M.G.} \\
      \dive{\vv{B}}=0 & \text{M.T.} 
    \end{cases}
    }
    % \end{equation}
  \item Les équations aux rotationnels (M.F. et M.A.) traduisent le couplage entre les champs $\E$ et $\B$,
  %  \begin{equation}
    \begin{center}
      \boxed{
        \begin{cases}
          \rot{\vv{E}}=-\frac{\partial \vv{B}}{\partial t} & \text{M.F.} \\
          \rot{\vv{B}}=\mu_0\varepsilon_0\frac{\partial \vv{E}}{\partial t} & \text{M.A.}
        \end{cases}
        }
    \end{center}
    % \end{equation}
\end{liste}
  \end{remarque}

\item Le milieu est \imp{illimité}. 
\begin{remarque}
  Le champ ne vérifie alors aucune condition aux limites particulières, si ce n'est qu'il doit rester fini à l'infini.
\end{remarque}
\end{enumerate}
On montre alors que le champ électromagnétique satisfait à l'équation de d'Alembert.

\begin{methode}[Établir l'équation des ondes électromagnétique]
  \item	Appliquer l'opérateur $\rot$ à l'équation de \imp{Maxwell-Faraday}.
  \item Simplifier avec l'identité $\rot{\left(\rot{\vv{E}}\right)}=\grad{\left(\dive{\vv{E}}\right)}-\Delta\vv{E}$ et l'équation de \imp{Maxwell-Gauss}.
  \item Simplifier avec l'équation de \imp{Maxwell-Ampère}, et identifier l'équation de d'Alembert.
\end{methode}

\begin{demo}
% Les équations de Maxwell se simplifient alors en
% \begin{equation}
% \begin{cases}
% \dive{\vv{E}}=0 & \text{M.G.} \\
% \dive{\vv{B}}=0 & \text{M.T.} \\
% \rot{\vv{E}}=-\frac{\partial \vv{B}}{\partial t} & \text{M.F.} \\
% \rot{\vv{B}}=\mu_0\varepsilon_0\frac{\partial \vv{E}}{\partial t} & \text{M.A.}
% \end{cases}
% \end{equation}

Nous partons de équations au rotationnel pour isoler un des champs, par exemple $\E$. Prenons le $\rot$ de M.F.,
%  Les équations aux rotationnels (M.F. et M.A.) traduisent le couplage entre $\E$ et $\B$. Nous partons de ces équations-là pour isoler $\E$, par exemple :
\begin{equation}
\begin{split}
\rot{\left(\rot{\vv{E}}\right)}&=\grad{\left(\dive{\vv{E}}\right)}-\Delta\vv{E}=-\Delta\vv{E}\\
&=\rot{\left(-\frac{\partial \vv{B}}{\partial t}\right)}=-\frac{\partial }{\partial t}\rot{\vv{B}}=-\mu_0\varepsilon_0\frac{\partial^2\vv{E}}{\partial t^2}.
\end{split}
\end{equation}
On voit qu'on retrouve une équation de d'Alembert
% \begin{equation}
\eqbox[]{
  \Delta\vv{E}-\mu_0 \epz \frac{\partial^2\vv{E}}{\partial t^2}=0.
  }
% \end{equation}
La célérité des ondes électromagnétiques vaut $c=\frac{1}{\sqrt{\varepsilon_0\mu_0}} = \SI{299792458}{m.s^{-1}}$. C'est précisément la célérité de la lumière ! C'est ainsi que Maxwell comprit en 1964 que la lumière n'est rien d'autre qu'une onde électromagnétique. On montre de la même manière en calculant $\rot{\left(\rot{\vv{B}}\right)}$ que $\vv{B}$ satisfait à la même équation de d'Alembert. La grandeur $c$ désigne donc la vitesse de l'onde électromagnétique dans le vide.
\end{demo}

%Sa valeur est fixée par définition\footnote{Cela définit exactement le mètre, car la seconde est définie à partir de la durée correspondant à \num{9192631770} périodes de l'atome de césium 133.} à

\begin{theorem}[Équation des ondes électromagnétiques]\label{prop:}
Le champ électromagnétique $(\E,\B)$ dans le \imp{vide} satisfait à l'équation de d'Alembert
$$
\Delta\vv{E}- \frac{1}{c^2} \frac{\partial^2\vv{E}}{\partial t^2}=0, \qquad \Delta\vv{B}-\frac{1}{c^2} \frac{\partial^2\vv{B}}{\partial t^2}=0,
$$
de célérité la célérité de la lumière dans le vide $c=\frac{1}{\sqrt{\varepsilon_0\mu_0}} = \SI{299792458}{m.s^{-1}}$.
\end{theorem}

\begin{exemple}{Phénomènes propagatifs et équation de d'Alembert}
  Nous avons vu maintenant un certain nombre de phénomènes propagatifs décrits par l'équation de d'Alembert. Ces phénomènes propagatifs font tous intervenir deux grandeurs, appelées grandeurs duales, dont les variations temporelles et spatiales sont couplés par une loi de la physique. Remplir le tableau d'analogie suivant.
\begin{center}
  \begin{tabular}{|c|c|c|c|}
    \hline
    Phénomène propagatif & Grandeurs duales & Équation de couplage & Équation de couplage \\
    \hline
    Onde électromagnétique & $\begin{cases} \text{champ électrique} & \E \\ \text{champ magnétique} & \B \end{cases}$ & $\rot{\E}=-\dr{\B}{t}$ & $\rot{\B}=\mu_0\epz\dr{\E}{t}$ \\
    \hline
    Onde mécanique corde & & &  \\[1cm]
    \hline
    Onde acoustique fluide &  &  &  \\[1cm]
    \hline
    
  \end{tabular}
\end{center}

\end{exemple}
\examplesol{
  \begin{center}
    \begin{tabular}{|c|c|c|c|}
      \hline
      Phénomène propagatif & Grandeurs duales & Équation de couplage & Équation de couplage \\
      \hline
      Onde électromagnétique & $\begin{cases} \text{champ électrique} & \E \\ \text{champ magnétique} & \B \end{cases}$ & M.F. $\rot{\E}=-\dr{\B}{t}$ & M.A. $\rot{\B}=\mu_0\epz\dr{\E}{t}$ \\
      \hline
      Onde mécanique corde & $\begin{cases} \text{vitesse} & v=\dr{y}{t} \\ \text{angle} & \alpha = \dr{y}{x} \end{cases}$ & PFD $\mu \dr{v}{t}= T \dr{\alpha}{x}$ & Géométrie  $\dr{\alpha}{t}=\dr{v}{x}$ \\
      \hline
      Onde acoustique fluide & $\begin{cases} \text{surpression} & p_1 \\ \text{vitesse} & v_1 \end{cases}$ & Euler $\rho_0 \dr{v_1}{t}=-\dr{p_1}{x}$ & Conservation $\rho_0 \chi_S \dr{p_1}{t} = - \dr{v_1}{x}$ \\
      \hline
      
    \end{tabular}
  \end{center}
}

\section{Ondes planes progressives harmoniques (OPPH)}
Toute solution de l'équation de d'Alembert peut s'écrire comme la superposition d'OPPH. Dans la suite du chapitre, nous étudierons les solutions OPPH, qui
%Pour analyser la structure de l'onde et les solutions de l'équation d'onde, il va être pratique de chercher les solutions sous la forme d'ondes planes progressives harmoniques (OPPH) donc
en notation complexe s'écrivent
\begin{equation}
\label{eq:E_B_convention_complexe}
 \begin{cases}
 \underline{\vv{E}}(M,t)=\underline{\vv{E}_0}\,\e^{i(\omega t-\vv{k}\cdot\vv{r})},\\
  \underline{\vv{B}}(M,t)=\underline{\vv{B}_0}\,\e^{i(\omega t-\vv{k}\cdot\vv{r})}.
 \end{cases}
 \end{equation}
Nous posons $\vv{k}=k\vv{u}$ pour le vecteur d'onde avec $\omega=kc$ d'après la relation de dispersion et où $\vv{u}$ est le vecteur unitaire qui donne la direction et le sens de propagation de l'onde.
% et qui est par définition orthogonal aux plans d'onde.

\subsection{Équations de Maxwell en notation complexe}


\parag{Opérateurs de dérivation en notation complexe}
Considérons l'action des opérateurs de dérivées temporelle et spatiale sur une OPPH se propageant dans la direction $\vv{u}$, par exemple sur le champ électrique 
$$
\udl{\vv{E}}(M,t)=\udl{\vv{E}_0}\,\e^{i(\omega t-k_x x - k_y y - k_z z)}.
$$
En dérivant par rapport au temps, il vient
$$
\dr{\udl{\E}}{t} = \udl{\Vec{E_0}}\times i \omega \e^{i(\omega t-k_x x - k_y y - k_z z)} = i\omega \udl{\E}.
$$
On en déduit que la dérivée temporelle agit en notation complexe comme
$$
\dr{}{t} = i\omega.
$$
De plus, en calculant une dérivée spatiale, on trouve
$$
\dr{\udl{\E}}{x} =  \udl{\Vec{E_0}}\times (-i k_x) \e^{i(\omega t-k_x x - k_y y - k_z z)} = -ik_x \udl{\Vec{E}}.
$$
On trouve une relation similaire pour les dérivées par rapport à $y$ et $z$.
On peut définir l'opérateur vectoriel \guill{nabla} par
$$
\Vec{\nabla} = \left(\dr{}{x},\dr{}{y},\dr{}{z} \right).
$$
En notation complexe, l'opérateur nabla agit comme
$$
\Vec{\nabla} = - i (k_x,k_y,k_z) = - i\Vec{k}.
$$


\begin{prop}[Opérateurs de dérivation en notation complexe]\label{prop:}
  Les opérateurs en notation complexe agissent sur une OPPH comme
  \begin{equation*}
  % \begin{cases}
  \frac{\partial}{\partial t} =  i\omega \qquad\text{ et }\qquad
  \vv{\nabla} = -i\vv{k}.
  % \end{cases}
  \end{equation*}
  On peut déterminer l'action des opérateurs usuels d'analyse vectorielle avec la correspondance
$$
\grad f= \Vec{\nabla} f = - if \Vec{k},  \quad \dive \Vec{v} = \Vec{\nabla} \ps \Vec{v} = - i \Vec{k} \ps \Vec{v}, \quad \rot \Vec{v} = \Vec{\nabla} \pv \Vec{v} = - i \Vec{k} \pv \Vec{v}.
$$
\end{prop}
% \vspace{-1cm}

% \begin{rappel}{Opérateurs vectoriels en notation complexe}
% On rappelle les équivalences suivantes pour les opérateurs en notation complexe,
% 	%et en remarquant qu'avec les conventions \eqref{eq:E_B_convention_complexe} nous avons l'équivalence suivante\footnote{L'équivalence pour les dérivées spatiales est simple à voir si on décompose terme à terme en coordonnées cartésiennes: $\vv{k}\cdot\vv{r}=k_{x}x+k_{y}y+k_z z$ et donc $\frac{\partial }{\partial x}=-ik_x$, $\frac{\partial }{\partial y}=-ik_y$, $\frac{\partial }{\partial z}=-ik_z$. Finalement en généralisant $\frac{\partial }{\partial x}\vv{e_x}+\frac{\partial }{\partial y}\vv{e_y}+\frac{\partial }{\partial z}\vv{e_z}=\vv{\nabla}=-i\vv{k}$.}
% 	\begin{equation}
% 	\begin{cases}
% 	\frac{\partial}{\partial t} ~\leadsto~ i\omega, \\[2mm]
% 	\vv{\nabla}~\leadsto~ -i\vv{k}.
% 	\end{cases}
% 	\end{equation}
% 	On rappelle également que pour tout champ vectoriel $\vv{A}$, on a $\dive{\vv{A}}=\vv{\nabla}\cdot\vv{A}$, $\rot{\vv{A}}=\vv{\nabla}\pv\vv{A}$.
% \end{rappel}
\parag{Équations de Maxwell}
Les équations de Maxwell se réécrivent donc en notation complexe comme
\begin{equation}
\begin{cases}
  -i\vv{k}\cdot\vv{\underline{E}}= 0,\\
- i\vv{k}\cdot\vv{\underline{B}}=0, 
\end{cases}\quad\text{et}\qquad
\begin{cases} 
-i\vv{k}\pv\vv{\underline{E}}=-i\omega\vv{\underline{B}},\\
-i\vv{k}\pv\vv{\underline{B}}=\mu_0 \epz i\omega\vv{\underline{E}}.
\end{cases}
\end{equation}
ce qui donne les relations
\eqbox[]{
  \begin{cases}
    \vv{k}\cdot\vv{\underline{E}}=0, \\
    \vv{k}\cdot\vv{\underline{B}}=0, 
    \end{cases}\quad\text{et}\qquad
    \begin{cases} 
      \vv{k}\pv\vv{\underline{E}}=\omega\vv{\underline{B}},\\
      \vv{k}\pv\vv{\underline{B}}=-\frac{\omega}{c^2}\vv{\underline{E}}.
    \end{cases}
}

\subsection{Structure de l'onde}

\parag{Caractère transverse des ondes électromagnétiques}
L'équation de M.T., qui est valable même dans un milieu matériel, indique que $\udl{\Vec{B}}$ est orthogonal au vecteur d'onde $\Vec{k}$, et donc à la direction de propagation de l'onde. L'équation de M.G. montre que $\udl{\Vec{E}}$ est aussi  orthogonal à la direction de propagation de l'onde dans le vide. Ces propriétés restent valables pour les champs réels !
\begin{prop}[Caractère transverse]\label{prop:}
		L'onde électromagnétique dans le \imp{vide} est une onde \important{transverse}, car $\E$ et $\B$ sont orthogonaux à $\Vec{k}$.
\end{prop}

\parag{Relation de structure} On cherche maintenant une relation entre les deux grandeurs duales $\E$ et $\B$. L'équation de M.F., qui est valable même dans un milieu matériel, donne la relation générale
% \begin{equation}
  \eqbox[]{
    \vv{\udl{B}}=\frac{\vv{k}\pv\vv{\udl{E}}}{\omega}
  }
  % \end{equation}
En utilisant la relation de dispersion $\vv{k}=\frac{\omega}{c}\vv{u}$, valable uniquement dans le vide, on trouve de plus
  \begin{equation}
    \vv{\udl{B}} =\frac{\vv{u}\pv\vv{\udl{E}}}{c}.
  \end{equation}
  Les champs $\vv{\udl{E}}$ et $\vv{\udl{B}}$ vibrent en phase, car cette relation de structure ne fait intervenir que des nombres réels. Elle est donc également valable pour les champs réels $\E$ et $\B$ !
  De plus, comme $\E$ est orthogonal à $\Vec{k}$ dans le vide, la famille $(\vv{u},\vv{E},\vv{B})$ forme un trièdre orthogonal direct.
  Par passage à la norme nous avons aussi la relation utile
$E=cB.$
L'équation de M.A. est automatiquement satisfaite si celle de M.F. et la relation de dispersion $\omega=ck$ le sont. Elle n'apporte pas davantage d'information, mais permet réciproquement de déterminer directement $\vv{E}$ lorsqu'on connaît $\vv{B}$, car
\begin{equation}
-\vv{k}\pv\vv{B}=\frac{\omega}{c^2}\vv{E}\quad\Rightarrow\quad \boxed{\vv{E}=-\vv{u}\pv c\vv{B}.}
\end{equation}

\begin{prop}[Relation de structure]\label{prop:}
	Le champ électromagnétique d'une OPPH dans le \imp{vide} satisfait la relation de structure
$$
\vv{B}=\frac{\vv{u}\pv\vv{E}}{c} \Leftrightarrow \vv{E}=-\vv{u}\pv c\vv{B}.
$$
Ainsi, $(\vv{u},\vv{E},\vv{B})$ forme un trièdre orthogonal direct et $E=cB$ en norme.
\end{prop}

% \begin{remarque}
% Attention à bien se rappeler que ces résultats ne sont valables que pour une OPPH solution de l'équation de d'Alembert dans le vide illimité, qui satisfait donc à la relation de dispersion $\omega=ck$.
% \end{remarque}



\begin{exemple}{Champ électromagnétique d'une OPPH dans le vide}
  On étudie une OPPH électromagnétique dans le vide.
  \begin{quest}
    \item	On donne le champ électrique $\vv{E}=E_0\,\cos(\omega t-kz)\ex$. Déterminer le champ magnétique.
    \item On donne le champ magnétique $\vv{B}=B_0\,\cos(\omega t-kz)\ey$. Déterminer le champ électrique.
  \end{quest}
\end{exemple}
\examplesol{
  \begin{quest}
    \item	Par exemple si on donne dans le vide $\vv{E}=E_0\,\cos(\omega t-kz)\vv{e_x}$, alors on en déduit immédiatement  que $\vv{u}$ est selon $\vv{e_z}$ (les plans d'onde correspondent à $z=\cte$ et sont orthogonaux à la direction de propagation fixe). Ainsi 
    \eqbox{
      \vv{B}=\frac{\vv{e_z}\pv\vv{E}}{c}=\frac{E_0}{c}\cos(\omega t-kz)\vv{e_y}. 
      }
      La seule connaissance du champ électrique suffit à déterminer le champ magnétique. On peut par exemple ensuite déterminer le vecteur de Poynting, la densité volumique d'énergie du champ, etc.
    \item De même d'après la forme de $\vv{B}$ la direction de propagation est selon $\vv{e_z}$, et on obtient alors que 
    \eqbox{
      \vv{E}=-\vv{e_z}\pv  c B_{0}\cos(\omega t-kz)\vv{e_y}=cB_0\cos(\omega t-kz)\vv{e_x}. 
      }
      On retrouve bien le cas précédent.

  \end{quest}
}


\subsection{Considérations énergétiques}

\begin{danger}
  
  Il faut toujours revenir aux \important{notations réelles} lorsqu'on calcule des \important{grandeurs énergétiques} ! En effet, il s'agit de grandeurs quadratiques, or
  %\footnote{Je déconseille aussi l'utilisation de $\frac{1}{2}\text{Re}\left(\vv{\underline{E}}\pv\vv{\underline{B}^{\star}}\right)$, car bien que correcte, il est rapide d'oublier dans cette relation le facteur $\frac{1}{2}$, ou pire ...}
  \begin{equation}
    \text{Re}\left(\vv{\underline{E}}\right)\pv\text{Re}\left(\vv{\underline{B}}\right)\JRouge{\boldsymbol{\neq}}~\text{Re}\left(\vv{\underline{E}}\pv\vv{\underline{B}}\right)
  \end{equation}
\end{danger}

\parag{Répartition d'énergie} Rappelons que l'énergie électromagnétique est répartie entre les champs électriques et magnétiques, suivant la densité volumique
$$
u_{\rm{em}}=u_{\rm e} + u_{\rm m} = \frac{1}{2}\varepsilon_0 E^2+\frac{B^2}{2\mu_0}.
$$
Dans le \imp{vide}, il n'y a pas de courant  électrique $\j$ et donc pas de perte volumique par effet Joule, soit $p_{\rm vol} = 0$. L'équation de conservation de l'énergie se simplifie alors en 
$$
\dr{u_{\rm em}}{t} + \dive{\vv{\Pi}} = 0,
$$
où $\Pi$ désigne le vecteur de Poynting.
Pour une \imp{OPPH} dans le vide, nous avons vu que $E=cB$. Les densités volumiques d'énergie valent donc 
\begin{equation}
u_{\rm{e}}=\frac{1}{2}\varepsilon_0 E^2 = \frac{1}{2}\varepsilon_0 c^2 B^2 =  \frac{B^2}{2\mu_0} = u_{\rm{m}}.
\end{equation}
Comme pour les ondes acoustiques, il y a \imp{équipartition} de la densité volumique d'énergie du champ, ici entre les formes électriques et magnétiques. La densité totale d'énergie électromagnétique vaut 
\eqbox[]{
  u_{\rm{em}}=\varepsilon_0 E^2.
}

\parag{Vecteur de Poynting}
On considère une OPPH donc d'après la relation de structure $\vv{B}=\frac{\vv{u}\pv\vv{E}}{c}$. En utilisant que $\vv{a}\pv(\vv{b}\pv\vv{c})=\vv{b}(\vv{a}\cdot\vv{c})-\vv{c}(\vv{a}\cdot\vv{b})$ et $\vv{u}\cdot\vv{E}=0$ (Maxwell-Gauss), le vecteur de Poynting s'exprime alors comme\footnote{Attention à \important{toujours} prendre des champs réels pour calculer les grandeurs énergétiques!}
\begin{equation}
\vv{\Pi}=\frac{\vv{E}\pv\vv{B}}{\mu_0}=\frac{\vv{E}\pv(\vv{u}\pv\vv{E})}{c\mu_0}=\frac{E^2}{c\mu_0}\vv{u}=c \varepsilon_0 E^2\vv{u}.
\end{equation}
 On obtient donc la relation
% \begin{equation}
  \eqbox{
    \Poy = c\, u_{\rm{em}} \Vec{u}.
  }
% \end{equation}
Cela rappelle la relation analogue obtenue pour les ondes acoustiques. L'interprétation physique est identique: $c$ donne la vitesse de propagation de l'énergie de l'onde électromagnétique, et le vecteur de Poynting dirige donc le flux énergétique de l'onde électromagnétique. Pour l'OPPH dans le vide, le vecteur de Poynting est colinéaire à la direction de propagation de l'onde.

\begin{prop}[Énergétique d'une OPPH électromagnétique dans le vide]\label{prop:}
  Il y a \imp{équipartition} entre les formes électriques et magnétiques de l'énergie électromagnétique d'une OPPH. L'énergie se déplace à la \imp{célérité $c$} de la lumière dans la direction de propagation $\Vec{u}$.
\end{prop}


\begin{exemple}{Lumière solaire}
  Le flux surfacique radiatif solaire incident sur la Terre vaut $\vphi_S = \SI{1.36e4}{W/m^2}$. 
  \begin{quest}
    \item	Calculer l'amplitude du champ électrique $E_0$ associé au flux solaire.
    \item Calculer l'amplitude du champ magnétique $B_0$ associé au flux solaire.
  \end{quest}
\end{exemple}
\examplesol{
  \begin{quest}
    \item	Par définition, le flux surfacique s'identifie à la moyenne temporelle de la norme du vecteur de Poynting. Pour une OPPH dans le vide, on a 
    $$
    \vphi_S = \moy{\norme{\Poy}} = c \epz \moy{E^2} = c \epz E_0^2 \moy{\cos^2(\omega t - k z)}, 
    $$
    soit 
    $$
    \vphi_S = \frac{1}{2} c \epz E_0^2,
    $$
    et donc 
    \eqbox{
      E_0 = \sqrt{\frac{2\vphi_S}{c \epz}} = \SI{3.20e3}{V.m^{-1}}.
    }
    \item Pour une OPPH dans le vide, on a 
    \eqbox{
      B_0 = \frac{E_0}{c}  = \SI{1.07e-5}{T}.
    }
  \end{quest}
}

\parag{Interprétation corpusculaire du rayonnement}
Qu'est ce qui véhicule l'énergie du champ électromagnétique ? À l'échelle microscopique, une théorie quantique du rayonnement est nécessaire. Le champ électromagnétique doit être quantifié et décrit par des particules élémentaires de masse nulle qui transportent l'énergie du champ: les photons. On associe à chaque photon une énergie et une quantité de mouvement (ou impulsion) donnée par les relations de Planck-Einstein et de De Broglie.

\begin{theorem}[Interprétation corpusculaire du rayonnement]
  Le rayonnement électromagnétique peut être vu comme un flux de particules élémentaires de masse nulle, les \imp{photons}, d'énergie $E$ et de quantité de mouvement $\vv{p}$ données par 
  $$\begin{cases}
    E=\hbar\omega\qquad&\text{(relation de Planck-Einstein),}\\
    \vv{p}=\hbar\vv{k}\qquad&\text{(relation de De Broglie).}
  \end{cases}$$
\end{theorem}
% % \begin{equation}
% \eqbox[photon]{
%   }
% % \end{equation}

\begin{exemple}{Flux de photons d'un laser}
Considérons un laser He-Ne de puissance $\mc{P}=\SI{1}{mW}$ émettant un rayonnement à \SI{633}{nm}. Déterminer le flux de photons $\Phi$ du faisceau laser, i.e. le nombre de photons traversant une section unité du faisceau par unité de temps.
\end{exemple}
\examplesol{
  Soit $\d N$ le nombre de photons traversant une section du faisceau pendant une durée $\d t$. On a par définition, 
$$
\Phi=\frac{\d N}{\d t}.
$$
Soit $\d E =E  \d N$ l'énergie totale portée par ces $\d N$ photons, d'énergie individuelle $E$. La puissance $\mc{P}$ rayonnée par le laser est donnée par 
$$
\mc{P}=\frac{\d E}{\d t} = E \frac{\d N}{\d t} = \frac{h c}{\lambda}  \Phi.
$$
On en déduit
% \begin{equation}
\eqbox{
  \Phi = \frac{\mc{P}\lambda}{hc} \simeq \num{3e15} \text{ photons par seconde}.
  }
% \end{equation}
 En général, le nombre de photons est tellement grand que le caractère corpusculaire du rayonnement électromagnétique n'est pas visible. La description en terme de champ continu dans tout l'espace et donc d'onde électromagnétique est alors valable.
}

\section{Polarisation}

\subsection{Définition}

Considérons une OPPH qui se déplace dans la direction $\Vec{u} = \ez$. L’onde étant transverse, le champ électrique a l’expression générale suivante,
%
% Nous venons de voir à la question précédente que toute la structure de l'OPPH dans le vide est déterminée à partir de celle du champ $\vv{E}$: sa norme et sa direction. La direction de $\vv{E}$ dans un plan d'onde est appelée la direction de polarisation de l'onde. De façon complètement générale, puisque l'onde électromagnétique est transverse, si on suppose qu'elle se propage selon $\vv{e_z}$ alors on peut écrire en toute généralité
 \begin{equation}\label{eq:polar}
 \vv{E}(z,t) =E_{x}(z,t)\vv{e_x}+E_{y}(z,t)\vv{e_y} = E_{0x}\cos(\omega t - kz )\vv{e_x} + E_{0y}\cos(\omega t - kz + \varphi)\vv{e_y},
 \end{equation}
 où $k$ est la norme du vecteur d'onde et $\omega = k c$ la pulsation.
 Dans un plan d’onde $z=\cte$, l’extrémité du vecteur $\E(z,t)$ décrit une courbe particulière. La nature de cette courbe définit la \imp{polarisation}. Pour décrire cette courbe et déterminer complètement l'onde électromagnétique, il faut préciser le déphasage $\vphi \in ]-\pi,\pi]$ et le rapport entre les amplitudes $E_{0x}$ et $E_{0y}$.

 \begin{definition}[Polarisation]\label{def:}
   La \imp{polarisation} d’une OPPH électromagnétique désigne la courbe que décrit l’extrémité de $\E$ dans un plan d’onde. Elle est déterminée par le déphasage $\vphi$ et l'amplitude relative entre $E_{0x}$ et $E_{0y}$.
 \end{definition}

 Pour simplifier, nous choisirons le plan d'onde $z=0$ dans la suite. La polarisation est donc une courbe dans le plan $(\ex,\ey)$. Par convention, la polarisation est toujours définie de sorte que l'observateur voie l'onde arriver \imp{vers lui}. Le vecteur $\ez$ point donc vers l’observateur, si bien que $(\ex,\ey)$ forme une base directe du plan d’onde. On distingue différents types de polarisation : rectiligne, circulaire, elliptique, non polarisé. On trouvera des animations sur la polarisation à \url{https://youtu.be/Q0qrU4nprB0}.

\subsection{Polarisation rectiligne}
Une polarisation est \imp{rectiligne} lorsque $\E$ décrite une \imp{droite} dans un plan d'onde. Cela signifie que les composantes $E_y(t)$ et $E_x(t)$ oscillent en phase. On peut définir un angle $\alpha \in [0,\pi]$ tel que 
\begin{equation}
  E_y(t) = \tan(\alpha) E_x(t).
\end{equation}
La grandeur algébrique $\tan \alpha$ représente la pente de la droite.
Comme $E_x(t) = E_{0x}\cos(\omega t)$, on peut poser $E_{0x} = E_0 \cos\alpha$, ce qui implique $E_{0y} = E_0 \sin\alpha$ et $\vphi = 0$ ou $\pi$ suivant le signe de $\tan \alpha$.
On a alors
\begin{equation}
  \vv{E} = E_{0x} \cos(\omega t)\vv{e_x} + E_{0y}\cos(\omega t)\vv{e_y} = E_0 \cos(\omega t) \cos\alpha \, \ex + E_0 \cos(\omega t) \sin \alpha \, \ey = E_0 \cos(\omega t) \Vec{\veps},
\end{equation}
où l'on a définit le \imp{vecteur polarisation} $\Vec{\veps} = \cos\alpha \, \ex + \sin \alpha \, \ey$. 


\begin{figure}[H]
  \centering
  \begin{subfig}[label][0.4]
    \includegraphics{rect-polarization.pdf}
  \end{subfig}
  \begin{subfig}[label][0.2]
    \includegraphics{rect_gauche.pdf}
  \end{subfig}
  \begin{subfig}[label][0.2]
    \includegraphics{rect_droite.pdf}
  \end{subfig}
  \caption{(a) Représentation spatiale d'une OPPH électromagnétique polarisée rectilignement.  (c) Polarisation rectiligne pour $\alpha \in [0,\frac{\pi}{2}[ $ ; le déphasage vaut $\vphi=0$ et les composantes oscillent en phase. (b) Polarisation rectiligne pour $\alpha \in [\frac{\pi}{2},\pi[$ ; le déphasage vaut  $\vphi=\pi$ et les composantes oscillent en opposition de phase.}
  \label{fig:}
 \end{figure}
 
 \begin{definition}[Polarisation rectiligne]\label{def:}
 Une polarisation est \imp{rectiligne} lorsque $\E$ décrit une \imp{droite} dans un plan d'onde. Le champ électrique 
 $$
\E = E_0 \cos(\omega t) \Vec{\veps}
 $$
 est de déphasage $\vphi = 0$ ou $\pi$ et porté par le \imp{vecteur polarisation} $\Vec{\veps} = \cos\alpha \, \ex + \sin \alpha \, \ey$.
 \end{definition}
 
 \subsection{Polarisation circulaire}
 Une polarisation est \imp{circulaire} lorsque $\E$ décrit un \imp{cercle} dans un plan d'onde.
 Les composantes $E_x(t)$ et $E_y(t)$ vérifient l'équation 
 \begin{equation}
  E_{x}(t)^2+E_{y}(t)^2=E_{0}^2,
\end{equation}
ce qui représente un cercle de rayon $E_0>0$. On peut poser $E_x(t) = E_{0x} \cos(\omega t ) = E_0 \cos(\omega t)$. On voit alors que $E_y(t) = \pm E_0 \sin(\omega t) = E_0 \cos(\omega t \pm \frac{\pi}{2})$. En comparant à l'\Eq{polar}, il vient $\vphi = \pm \frac{\pi}{2}$.
\begin{liste}
  \item Lorsque $\vphi = + \frac{\pi}{2}$, le vecteur $\E$ tourne dans le sens indirect : la polarisation est circulaire \imp{droite}.
  \item Lorsque $\vphi = - \frac{\pi}{2}$, le vecteur $\E$ tourne dans le sens direct : la polarisation est circulaire \imp{gauche}.
\end{liste}
% Il s'agit d'un cas particulier de la polarisation elliptique, pour lequel $\varphi=\frac{\pi}{2}$ et $E_{0x}=E_{0y}=E_{0}$. La courbe décrite au cours du temps par le vecteur $\vv{E}$ dans un plan d'onde a alors pour équation
% On parle donc de polarisation circulaire.

En notation complexe, on a alors
\begin{equation}
  \udl{\E}(t) = E_0 e^{i\omega t} \ex + E_0 e^{i\omega t \pm \frac{\pi}{2}}\ey = E_0 e^{i\omega t}\left( \ex \pm i \ey \right).
\end{equation}
On peut définir formellement un vecteur de polarisation complexe
\begin{equation}
  \udl{\Vec{\veps}} = \ex \pm i \ey,
\end{equation}
où le signe $+$ s'applique à la polarisation circulaire droite, et le signe $-$ à une polarisation circulaire gauche. On a alors
\begin{equation}
  \udl{\E}(t) = E_0 e^{i\omega t} \udl{\Vec{\veps}}.
\end{equation}
\begin{remarque}
  Attention avec cette convention $\udl{\Vec{\veps}}$ n'est pas unitaire mais de norme $2$, donc la norme de $\E$ vaut
  \begin{equation}
    \E^2 = \frac{1}{2} \udl{\E} \ps \udl{\E}^* = \frac{E_0^2}{2} \udl{\Vec{\veps}} \ps \udl{\Vec{\veps}}^* = E_0^2.
  \end{equation}
\end{remarque}

\begin{definition}[Polarisation circulaire]\label{def:}
 Une polarisation est \imp{circulaire} lorsque $\E$ décrit un \imp{cercle} dans le plan d'onde. Le déphasage vaut $\vphi = +\frac{\pi}{2}$ (polarisation droite) ou $\vphi = -\frac{\pi}{2}$ (polarisation gauche). En notation complexe, on a
  \begin{equation}
    \udl{\E}(t) = E_0 \,e^{i\omega t} \Vec{\udl{\veps}},
  \end{equation}
  où le vecteur de polarisation complexe vaut $\Vec{\udl{\veps}} = \ex \pm i \ey$.
\end{definition}

\begin{figure}[H]
  \centering
  \begin{subfig}[label][0.4]
    \includegraphics{cyclic-polarization.pdf}
  \end{subfig}
  \begin{subfig}[label][0.2]
     \includegraphics{circu_droite.pdf}
   \end{subfig}
   \begin{subfig}[label][0.2]
     \includegraphics{circu_gauche.pdf}
   \end{subfig}
   \caption{(a) Représentation spatiale d'une OPPH électromagnétique polarisée circulairement. (b) Polarisation circulaire droite. (c) Polarisation circulaire gauche. }
   \label{fig:}
 \end{figure}


 \subsection{Polarisation elliptique}

Dans toutes les autres situations, i.e. lorsque le déphasage $\varphi \neq 0\, \text{mod}\, \frac{\pi}{2}$ est une constante quelconque, ou que $E_{0x} \neq E_{0y}$, nous avons pour $z=0$,
 \begin{equation}
 \begin{split}
 \vv{E}&=E_{0x}\cos(\omega t)\vv{e_x}+E_{0y}\cos(\omega t+\varphi)\vv{e_y}\\
 %&=E_{0x}\cos(\omega t)\vv{e_x}+E_{0y}\vv{e_y}.
 \end{split}
 \end{equation}

% Le demi-axe s'obtient en regardant qui de $\cos \alpha$ ou de $\sin \alpha$ est le plus grand. L'excentricité de l'ellipse vaut
% \begin{equation}
%   e = \frac{\sqrt{a^2-b^2}}{a} = \frac{\sqrt{\cos^2\alpha - \sin^2\alpha}}{\cos\alpha}
% \end{equation}
Dans le cas général où $\vphi$ est quelconque, les axes de l'ellipse sort tournés d'un certain angle par rapport aux axes $(Ox)$ et $(Oy)$.

\begin{definition}[Polarisation elliptique]\label{def:}
  Une polarisation est \imp{elliptique} lorsque $\E$ décrit une \imp{ellipse} dans le plan d'onde : les amplitudes $E_{0x}$ et $E_{0y}$ ou le déphasage $\vphi$ sont quelconques. L'ellipse est droite si $\vphi>0$ et gauche si $\vphi<0$.
\end{definition}

  % Dans le plan d'onde, le vecteur $\vv{E}$ décrit une ellipse. On dit que la polarisation est elliptique: au cours du temps l'extrémité du vecteur $\vv{E}$ vibre le long de cette ellipse, dont les axes diffèrent en général de $\vv{e_x}$ et $\vv{e_y}$. Le détail du calcul est proposé dans la remarque ci-dessous.

%
%  \begin{remarque}
% Pour déterminer l'équation de l'ellipse, nous allons écrire que $\cos^2(\omega t)+\sin^{2}(\omega t)=1$ en remplaçant $\cos(\omega t)$ et $\sin(\omega t)$ par leurs expressions. Tout d'abord, $\frac{E_x(t)}{E_{0x}}=\cos(\omega t)$, puis $\frac{E_y(t)}{E_{0y}}=\cos(\omega t+\varphi)$ qui se réécrit en développant comme
%  \begin{equation}
%  \frac{E_y(t)}{E_{0y}}=\left[\cos(\omega t)\cos(\varphi)-\sin(\omega t)\sin(\varphi)\right]=\left[\frac{E_x(t)}{E_{0x}}\cos(\varphi)-\sin(\omega t)\sin(\varphi)\right].
%  \end{equation}
%  En réorganisant pour isoler $\sin(\omega t)$ et d'après $\cos^2(\omega t)+\sin^{2}(\omega t)=1$ il vient
%  \begin{equation}
%  \left(\frac{E_x(t)}{E_{0x}}\right)^2+\left(\frac{E_x(t)\cos(\varphi)}{E_{0x}\sin(\varphi)}-\frac{E_y(t)}{E_{0y}\sin(\varphi)}\right)^2=1.
%  \end{equation}
%  Il s'agit effectivement de l'équation d'une ellipse dans un plan d'onde $(\vv{e_x},\vv{e_y})$. Les axes de l'ellipse sont confondus avec $(\vv{e_x},\vv{e_y})$ si et seulement si $\cos(\varphi)=0$ autrement dit pour $\varphi=\pm\frac{\pi}{2}$.
%  \end{remarque}
%  Enfin, rappelons que pour déterminer le sens de rotation de l'ellipse, on se place à deux instants bien choisis, usuellement $t=0$ et $t=T/4$, c'est-à-dire $\omega t=0$ et $\omega t=\frac{\pi}{2}$. L'hélicité de la polarisation est gauche si l'ellipse est parcourue dans le sens trigonométrique, et droite si l'ellipse est parcourue dans le sens horaire (avec toujours pour convention que l'onde arrive vers nous).

 \begin{figure}[H]
   \centering
   \begin{subfig}[label][0.4]
     \includegraphics{elliptical-polarization.pdf}
   \end{subfig}
   \begin{subfig}[label][0.2]
     \includegraphics{ellipse_droite.pdf}
   \end{subfig}
   \begin{subfig}[label][0.2]
     \includegraphics{ellipse_gauche.pdf}
   \end{subfig}
   \caption{(a) Représentation spatiale d'une OPPH électromagnétique polarisée elliptiquement. (b) Polarisation elliptique droite. (c) Polarisation elliptique gauche. }
   \label{fig:}
 \end{figure}

%  \begin{exemple}{Polarisation elliptiques}
%   On considère une polarisation elliptique $\vv{E}=E_{0x}\cos(\omega t)\vv{e_x}+E_{0y}\cos(\omega t+\varphi)\vv{e_y}$.
%   \begin{quest}
%     \item On suppose les amplitudes $E_{0x}$ et $E_{0y}$ différentes, et le déphasage égal à $\vphi=\pm\frac{\pi}{2}$. Montrer que les axes de l'ellipse sont confondus avec $(Ox)$ et $(Oy)$. Que vaut l'excentricité ?
%     \item	On suppose que les amplitudes $E_{0x} = E_{0y} =E_0$ sont égales. Pour quelle valeur de $\omega t$ la norme de $\E$ est-elle maximale ? minimale ? En déduire les axes de l'ellipse.
%   \end{quest}
%  \end{exemple}
%  \examplesol{
%   \begin{quest}
%     \item Considérons par exemple le cas où $E_{0x} > E_{0y}$ avec $\vphi=\frac{\pi}{2}$. Posons $E_{0x} = E_0$ et $E_{0y} = E_0 \sqrt{1-e^2}$ avec $0<e<1$. On trouve alors $E_x(t)=E_{0}\cos(\omega t)$ et $\frac{E_y(t)}{\sqrt{1-e^2}}=E_0 \sin(\omega t)$, d'où l'on tire
%     \begin{equation*}
%       E_x(t)^2 + \frac{E_y(t)^2}{1-e^2} = E_0^2.
%     \end{equation*}
%     Il s'agit de l'équation d'une ellipse d'axes $(Ox)$ et $(Oy)$, de demi-grand axe $E_{0x}=E_0$, de demi-petit axe $E_{0y}= E_0\sqrt{1-e^2}$ et d'excentricité $e$.
%     \item	La norme carré de $\E$ vaut 
%     \begin{equation*}
%       E^2 = E_{0x}^2\cos^2(\omega t) + E_{0y}^2\cos^2(\omega t+\varphi) = E_0^2\left(\cos^2(\omega t) + \cos^2(\omega t+\varphi)\right) = 
%     \end{equation*}
% \end{quest}
%  }


\subsection{Lumière non polarisée}

La lumière solaire n'entre dans aucune de ces catégories. En effet, lorsque le déphasage $\vphi(t)$ dépend du temps de façon aléatoire, et peut prendre toutes les valeurs possibles avec la même probabilité, la lumière explore toutes les polarisations possibles en un temps typique $\tau$ très petit devant le temps de récepteur $\tau_r$ des détecteurs usuels. La lumière apparaît donc, après moyenne temporelle, \imp{non polarisée}.
\begin{definition}[Lumière non polarisée]\label{def:}
  La phase $\vphi(t)$ d’une \imp{lumière non polarisée} est une \imp{fonction aléatoire} du temps.
\end{definition}


\end{document}
